\subsection{Response Fit, Severity and Source Identity}
\label{sec:narrative-information}

A better physical prediction need not produce a better location. Competing
sources may remain similar after their severities and initial conditions
are fitted. For a response $s_h$ at specified severity and context, the
distance to an admitted competitor is
\begin{equation}
 d_h^2=\min_{j\ne h}\ \min_{\beta_j\in\mathcal A_j}
       \|s_h-s_j(\beta_j)\|_{\Sigma^{-1}}^2.
 \label{eq:narrative-separation}
\end{equation}
The set $\mathcal A_j$ fixes allowed signs, severities and nuisance
coordinates. Small distance makes source discrimination noise-sensitive;
it is not a universal error formula or a calibrated decision certificate
\cite{eriksson2013diagnosability,scharf1994matched}.

Three outcomes remain separate: prediction error at a specified physical
source, ranking of competing sources, and uncertainty calibration. A
response improvement establishes the first. It establishes the second only
if source decisions improve under the same observations, horizon and
decision rule. Neither establishes the third automatically. A nonlinear
response comparison is needed before treating a surrogate's residual
confusion as an unavoidable physical limit.

The supplementary archive retains the earlier state-reconstruction and
simultaneous-load studies, with their separate observation contracts and
prior-normalization qualification. Their known-onset scores are not
combined with the causal results here. In particular, recovering hidden
voltages does not certify the source that moved them.

To first order, a component of the context derivative parallel to a nominal
source signature can be absorbed by fitting a different severity. Correct
source ranking can therefore coexist with biased intensity. Components
orthogonal to that signature cannot be removed by this rescaling. This
explains why source accuracy can reach a test ceiling while response and
severity errors still distinguish the models; it is a standard projection
argument, not a new identifiability theorem.
